% sec9_6.tex — Section 9.6: Application: Purchasing Power Parity

The theory of \textbf{purchasing power parity} (PPP) states that, for any two countries,
the currency exchange rate equals the ratio of the price levels of the two countries.
Put differently, once converted to a common currency, national price levels should
be equal. Let $P_t$ be the price index for the United States, $P_t^*$ be the price index for
the foreign country in question (say, the United Kingdom), and $S_t$ be the exchange
rate in dollars per unit of the foreign currency. PPP states that
\begin{equation}
\label{eq:9.6.1}
P_t = S_t \cdot P_t^*,
\end{equation}
or taking logs and using lower-case letters for them,\footnote{Equation~\eqref{eq:9.6.1}
or~\eqref{eq:9.6.2} is a statement of ``absolute PPP.'' What is called relative PPP requires only that
the first difference of~\eqref{eq:9.6.2} hold, that is, the rate of growth of the exchange rate should offset the differential
between the rate of growth in home and foreign price indexes. Relative PPP should not be confused with the
weaker version of PPP to be introduced below.}
\begin{equation}
\label{eq:9.6.2}
p_t = s_t + p_t^*.
\end{equation}
PPP is related to but different from the \textbf{law of one price}, which states that~\eqref{eq:9.6.1}
holds for any good, with $P_t$ representing the dollar price of the good in question
and $P_t^*$ the foreign currency price of the good. If international price arbitrage
works smoothly, the law of one price should hold at every date $t$. If the law of one
price holds for all goods and if the basket of goods for the price index is the same
between two countries, then PPP should hold.

Due to a variety of factors limiting international price arbitrage, PPP does not
hold exactly at every date $t$.\footnote{See, e.g., Section~4 of Rogoff (1996) for a list of factors that prevent PPP from holding.} A weaker version of PPP is that the deviation
from PPP
\begin{equation}
\label{eq:9.6.3}
z_t \equiv s_t - p_t + p_t^*,
\end{equation}
which is the real exchange rate, is stationary. Even if it does not enforce the law
of one price in the short run, international arbitrage should have an effect in the
long run.

\subsection*{The Embarrassing Resiliency of the Random Walk Model?}

For many years researchers found it difficult to reject the hypothesis that real
exchange rates follow a random walk under floating exchange regimes. As Rogoff
(1996) notes, this was somewhat of an embarrassment because every reasonable
theoretical model suggests the weaker version of PPP. However, studies since the
mid 1980s looking at longer spans of data have been able to reject the random walk
null. Recent work by Lothian and Taylor (1996) is a good example. It uses annual
data spanning two centuries from 1791 to 1990 for dollar/sterling and franc/sterling
real exchange rates to find that the random walk null can be rejected. In what follows, we apply the ADF $t$-test to their data on the dollar/sterling (dollars per pound)
exchange rate.

The dollar/sterling real exchange rate, calculated as~\eqref{eq:9.6.3} and plotted in Figure~9.4, exhibits no time trend. We therefore do not include time in the augmented
autoregression. If we were dealing with the weak version of the law of one price,
with $P_t$ and $P_t^*$ denoting the home and foreign prices of a particular good, then we
could exclude the intercept from the augmented autoregression because a change
in units of measuring the good does not affect $z_t$. But $P_t$ and $P_t^*$ here are price
indexes, and a change in the base year leads to adding a constant to $z_t$. To make
the test invariant to such changes, a constant is included in the augmented autoregression. By applying the ADF test, rather than the DF test, we can allow for serial
correlation in the first differences of $z_t$ under the null, but the lag chosen by
BIC (with $p_{\max}(T) = \big[12 \cdot (200/100)^{1/4}\big] = 14$) turned out to be 0, so the ADF $t$-test reduces to the DF $t$-test and an augmented autoregression reduces to an AR(1)
equation. The estimated AR(1) equation is
\[
z_t = \underset{(0.052)}{0.179} + \underset{(0.0326)}{0.8869}\; z_{t-1},
\quad SER = 0.071, \quad R^2 = 0.790, \quad t = 1792, \ldots, 1990.
\]

\begin{figure}[htbp]
\centering
\framebox[\textwidth]{\rule{0pt}{2.5in}\parbox{0.8\textwidth}{\centering
  Figure 9.4: Dollar/Sterling Real Exchange Rate, 1791--1990, 1914 Value Set to 0}}
\caption{Dollar/Sterling Real Exchange Rate, 1791--1990, 1914 Value Set to 0}
\label{fig:9.4}
\end{figure}

The DF $t$-statistic ($t^\mu$) is $-3.5$ ($= (0.8869 - 1)/0.0326$), which is significant at
1 percent. Thus, we can indeed reject the hypothesis that the sterling/dollar real
exchange rate is a random walk.

An obvious caveat to this result is that the sample period of 1791--1990 includes
fixed and floating rate periods. It might be that international price arbitrage works
more effectively under fixed exchange rate regimes. If so, the $z_{t-1}$ coefficient
should be closer to 0 during fixed rate periods. Lothian and Taylor (1996) report
that if one uses a simple Chow test, one cannot reject the hypothesis that the $z_{t-1}$
coefficient is the same before and after floating began in 1973. In the empirical
exercise of this chapter, you will be asked to verify their finding and also use the
ADF-GLS to test the same hypothesis.
